% =====================================================================
\section{v6.3: vNext+ Full Integration of Adaptive Tuning and StrongEngine $\Omega$ Full}
\label{sec:vnext-plus-v63}


\nrmobox{v7 structural correction}{
This mechanism resides on the \textbf{real side}: the StrongEngine $\Omega$
Full advances over each domain's true dynamics with a maximum-forward default,
while a civ-sim MaxForwardEngine imitates it to explore extremes. Rollout must
use the target domain's true dynamics with a long-enough horizon to credit
compounding, avoiding over-defensive collapse. Maximum forward is the default,
not a universal optimum (world-dependence). See Realignment Master Spec.
}
Sections~\ref{sec:stage11-v51}--\ref{sec:calibration-v62} described the
v5.1--v6.2 simulator development, in which the \texttt{Adaptive\_OmegaFull}
agent strategy was implemented as a simplified mutation~+~drift-penalty
heuristic rather than the full Adaptive NRMOvNext~+~StrongEngine~$\Omega$
Full pipeline of the Part~IX specification. v6.3 closes this gap.

\subsection{Discovery: implementation already existed}

The v52 codebase (\texttt{v52\_codebase/}, accompanying this monograph)
already contains the full implementation of:
\begin{itemize}
\item \texttt{governance/nrmo\_core.py} (92 lines): NRMO Core veto and
  admissibility-set construction, both original and vNext variants;
\item \texttt{governance/tuning\_layer.py} (83 lines): five-mode
  \texttt{MetaController} with hysteresis, \texttt{adaptive\_tuning}
  state-responsive rules, world-profile selection;
\item \texttt{engine/omega\_full.py} (703 lines): $\Omega$ Full's nine
  sections —~Candidate Population System (\S1), Favorable State Detector
  (\S4: Wolf Pursuit), Edge Survival Guard (\S6), Risk-Adjusted floor
  (\S7), Portfolio Synergy (\S8), Dual Objective Scoring (\S9), plus
  Long-Horizon Drift Estimator and Failure Memory.
\end{itemize}

The reason v5.0--v6.2 World Simulation did not use these was
architectural: v52 operates on a single \texttt{CivState} (six scalars:
$R,E,G,O,K,X$) whereas World Simulation operates on $N$~agents $\times$~6
per-agent dimensions. The two scales had not been bridged.

\subsection{Bridging design}

The v6.3 module \texttt{vnext\_plus.py} introduces three bridges:

\paragraph{Aggregation bridge.}
A per-civilisation \texttt{aggregate\_to\_civstate} function maps
$\{$agent population state arrays, religion strength, shock add$\} \to
\texttt{CivState}$. The mapping is heuristic but consistent with v52's
$0$--$130$ scaling:
\begin{equation*}
\begin{aligned}
R &= 130 \cdot \overline{\text{assets}} \\
K &= 130 \cdot \overline{\text{edu}} \\
G &= 130 \cdot \overline{\text{inst}} \\
O &= 130 \cdot (0.6\,\overline{\text{urban}} + 0.4\,\overline{\text{edu}})
    + 30\,\mathrm{Var}(\text{edu}+\text{assets}) \\
E &= \max(8, 65 - 80\,\text{shock} + 50\,\overline{\text{fk}}) \\
X &= \min(95, 100\,\text{shock} + 15\,\mathrm{Var}(\text{edu}+\text{assets}))
\end{aligned}
\end{equation*}

The inequality term $\mathrm{Var}(\text{edu}+\text{assets})$ in $O$ and
$X$ captures the previously ignored dimension: heterogeneity within the
civilisation contributes to both optionality (diverse positions $\Rightarrow$
more pathways) and exposure (inequality stress).

\paragraph{Per-civ controller.}
The \texttt{VNextPlusCivController} class instantiates one
\texttt{MetaController}, one \texttt{HysteresisTracker}, and one
\texttt{OmegaFullEngine} per civilisation. This is the principled
extension of v52's single-civilisation assumption to the World Simulation's
9-civilisation reality. Each controller maintains its own failure memory,
previous state, and world-archetype scores. Cross-civilisation
contamination of governance state is structurally precluded.

\paragraph{Action projection bridge.}
The civ-level optimal action $a^\star_c$ returned by
\texttt{OmegaFullEngine.select\_action} is projected onto agents via
\texttt{project\_action\_to\_agents}: for each agent of strategy $s$,
the agent's action is $a^\star_c + \mathcal{N}(0, \sigma_s^2 \mathbf{I}_4)$,
clipped and renormalised. The $\sigma_s$ values reflect strategy
fidelity to the civ optimum:
\texttt{NRMO\_vNext} = 0.03, \texttt{Adaptive\_OmegaFull} = 0.04,
\texttt{RiskAdj} = 0.06, \texttt{EVMax} = 0.05 (with $+0.05$ growth
bias), \texttt{Faith\_*} = 0.06--0.10 with sub-faith-specific biases,
\texttt{Drift} = $\infty$ (ignores civ action). Thus only the two
``vNext-family'' strategies tightly track the civ optimum; other
strategies use their original per-agent theory functions.

\subsection{Evolution beyond v52: archetype classifier enabled}

v52's $\Omega$ Full implementation contains a world-archetype classifier
(\texttt{classify\_world\_archetype}: vulnerable, race, stress,
stagnation, normal probabilities) and \texttt{get\_edge\_profile} that
would adapt Edge Guard portfolio weights to detected archetype. In v52,
both were \emph{computed but not used} (the inline comment reads:
``computed, logged; disabled until validated at 500-run validation'').

v6.3 re-enables both. The 500-run validation gate is replaced by the
empirical multi-civilisation, multi-seed averaging that v6.3 itself
provides: nine civilisations $\times$ three seeds = 27 independent
trajectories per benchmark, sufficient to detect systematic effects.
This is not a deferral of the original validation requirement but a
substitution of a denser empirical regime.

\subsection{Pipeline summary}

For each civilisation $c$ at each generation $g$:
\begin{enumerate}
\item Aggregate agent population $\to$ \texttt{CivState}$_c^g$.
\item World parameters $\mathit{wp}_c$ derived from Cultural Module's
  era $i$ at year $y = 40g$ (severity mapping
  $\mathrm{base\_failure} \mapsto \{\text{shock\_prob, tail\_prob,
  rivalry, env\_drag, \ldots}\}$).
\item \texttt{MetaController}$_c$.update($s, \mathit{wp}_c$) $\to$
  mode~$m_c \in \{$Normal, HighStakes, Recovery, StagnationEscape,
  Race$\}$, with hysteresis (enter\_th $=2$, exit\_th $=3$).
\item \texttt{adaptive\_tuning}(profile$(\mathit{world\_name})$,~$s$,~$m_c$)
  $\to$ tuned \texttt{TuningConfig} $\mathit{tc}_c$.
\item Candidate pool: \texttt{build\_candidate\_pool}$(s, \mathit{wp}_c)$
  generates $\sim$23--29 candidates from base templates, mutation,
  synthesis, and invention; Wolf Pursuit expands the pool when
  favourable state is detected.
\item NRMO Core: \texttt{construct\_admissible\_set}(pool,~$s$,~mode=vnext,
  $\mathit{tc}_c$) returns admissible subset.
\item $\Omega$ Full select: for each admissible candidate, MC rollout
  ($\mathrm{depth}=6$, $\mathrm{repeats}=6$; doubled for Wolf), score
  with full nine-component objective; rank by score; build portfolio
  (main + hedge + probe) with Synergy bonus and drift-safe hedge
  selection.
\item Project $a^\star_c$ onto agents of strategy
  \texttt{NRMO\_vNext} and \texttt{Adaptive\_OmegaFull}; other
  strategies use original per-agent theory functions.
\end{enumerate}

Architectural invariants are preserved without exception: NRMO Core only
issues YES/NO/HOLD, $\Omega$ Full searches only inside the admissible
set, Tuning Layer adjusts thresholds but does not enter NRMO scoring.

\subsection{Empirical comparison}

Three-seed averaged comparison, scale~=~small (9~civilisations $\times$
5{,}000 agents, 75 generations), with vs without vNext+ enabled:

\begin{center}
\small
\begin{tabular}{llrrr}
\toprule
Strategy & World & Baseline & vNext+ & $\Delta$ \\
\midrule
\multirow{4}{*}{NRMO\_vNext}
  & Japan       & 2.148 & 2.146 & $-0.002$ \\
  & China       & 2.109 & 2.140 & $+0.031$ \\
  & Europe      & 0.769 & 0.807 & $+0.038$ \\
  & Polynesian  & 0.943 & 1.065 & $+0.122$ \\
\midrule
\multirow{4}{*}{Adaptive\_OmegaFull}
  & Japan       & 2.175 & 2.137 & $-0.038$ \\
  & Europe      & 0.820 & 0.803 & $-0.017$ \\
  & Polynesian  & 0.947 & 1.042 & $+0.094$ \\
  & IndigenousAmericas & 0.248 & 0.258 & $+0.011$ \\
\midrule
\multirow{2}{*}{ExpectedValueMax}
  & Europe      & 0.950 & 0.981 & $+0.032$ \\
  & IndigenousAmericas & 0.198 & 0.236 & $+0.038$ \\
\bottomrule
\end{tabular}
\end{center}

The pattern is asymmetric and matches the theoretical prediction:
vNext+ \emph{improves performance in hostile worlds} (Polynesian,
Europe, IndigenousAmericas) for vNext-family strategies, but
\emph{degrades performance in peaceful worlds} (Indic, Islamic) where
the Adaptive Tuning Layer's protective threshold tightening costs
otherwise-available growth opportunity. This is the classical
\emph{insurance premium} of vNext: pay efficiency in calm times to
secure survival under stress.

In peaceful worlds, \texttt{ExpectedValueMax} still wins on score, as
predicted by Q4 of the original vNext specification (Section~XXX
research questions). vNext+ is not a free win; it is the systematic
trading of peacetime efficiency for hostile-world resilience.

\subsection{Runtime cost}

Adding vNext+ approximately quadruples small-scale runtime (small
preset: 4\,s $\to$ 17\,s, $\times 4.3$) due to the per-civilisation MC
rollouts ($\sim$14 candidates $\times$ 6 rollouts $\times$ 6 depth per
civ per generation $\approx$ 500 transitions per civ per generation;
$\times$ 9 civs $\times$ 75 generations $\approx$ 340{,}000 additional
transitions per run). The cost is acceptable for small and medium
scale; large scale (1.8\,M agents) takes proportionally longer.

\subsection{Limits and outlook}

The aggregation bridge necessarily loses information: a civilisation
where half the agents are highly educated and half are not, and one
where all agents have median education, both produce the same
$\overline{\text{edu}}$ and therefore the same $K$. The variance terms
in $O$ and $X$ partially address this. A future $v$ might bridge
\emph{distributional} rather than \emph{mean} statistics into
CivState, at the cost of expanding the state-vector dimension beyond
the six-scalar v52 form.

The action projection assumes that agents within a strategy class are
exchangeable. This is true for the simulation purposes here, but
agents differing in initial occupation, region, or family lineage may
in principle warrant differentiated actions. Future iterations could
condition $\sigma_s$ on per-agent state.

Stage~4's GPU acceleration remains environment-blocked; vNext+ at
large scale ($n \ge 200{,}000$ per civ) takes minutes-scale runtime
which is acceptable for empirical work but precludes the $10^7$ and
$10^8$ Stage~3/4 targets.
